\section{Hardware keys}

Hardware and tree verbs are top-level commands (\code{generate},
\code{split}, \code{bind}, and so on).

\begin{lstlisting}
# Generate a keypair. The private key is printed to stdout ONCE and never
# written to disk by this tool -- redirect it yourself.
keyquorum generate --type encryption --public-key-out alice.pub \
  --label alice --register > alice.key
keyquorum register --type encryption --label bob --public-key-file bob.pub
keyquorum list
\end{lstlisting}

\code{generate --register} writes the public key and records it in one
step. \code{register} brings in a public key that already exists elsewhere
(a PEM or OpenSSH file, or another person's \file{.pub}) under a chosen
label. \code{list} shows every registered hardware key and every live split
tree.

\begin{noteBox}
Two key types are used throughout this manual: \flag{--type encryption}
keys unwrap sealed shares and envelopes; \flag{--type signing} keys produce
and verify Ed25519 signatures (private bridges, \code{sign} / \code{verify},
authenticated updates). A label commonly needs both --- an encryption key
to receive things, and a signing key to authorize or countersign them.
\end{noteBox}
